% 主小行星带（综述）
% license CCBYSA3
% type Wiki

本文根据 CC-BY-SA 协议转载翻译自维基百科\href{https://en.wikipedia.org/wiki/Asteroid_belt}{相关文章}。

\begin{figure}[ht]
\centering
\includegraphics[width=6cm]{./figures/f9c18c4daeef0b22.png}
\caption{} \label{fig_belt_1}
\end{figure}
\begin{figure}[ht]
\centering
\includegraphics[width=6cm]{./figures/cf13395dc35061d2.png}
\caption{到目前为止，腰带内最大的物体是矮行星谷神星。小行星带的总质量明显小于冥王星大约是冥王星月亮的两倍卡隆。} \label{fig_belt_2}
\end{figure}
小行星带是太阳系中的一个圆环形区域，以太阳为中心，大致横跨行星木星和火星轨道之间的空间。该区域包含大量固态、不规则形状的天体，称为小行星或小行星状星体。这些已被确认的天体大小各异，但都远小于行星，且平均彼此相距约一百万千米（约六十万英里）。这个小行星带也被称为主小行星带或主带，以区别于太阳系中其他小行星族群$^{[1]}$。

小行星带是太阳系中最小且最内侧的环星盘。位于其他区域的小型太阳系天体类别，还包括近地天体、半人马天体、柯伊伯带天体、散射盘天体、塞德纳型天体以及奥尔特云天体。主带总质量中约 60\% 集中在四颗最大的小行星：谷神星（Ceres）、灶神星（Vesta）、智神星（Pallas）和健神星（Hygiea）中。小行星带的总质量估计约为月球质量的 3\%$^{[2]}$。

谷神星是小行星带中唯一足够大以被归类为矮行星的天体，其直径约为 950 km；而灶神星、智神星和健神星的平均直径均小于 600 km$^{[3][4][5][6]}$。其余在矿物学上已被分类的天体尺寸一直延续到仅数米大小$^{[7]}$。小行星带中的物质分布极为稀疏，以至于多艘无人航天器已多次穿越该区域而未发生任何碰撞事故$^{[8]}$。尽管如此，大型小行星之间的碰撞仍然会发生，并能形成小行星族群，其成员具有相似的轨道特征和成分。带内各个小行星依据其光谱进行分类，大多数可归入三大基本类型：富碳的 C 型、富硅酸盐的 S 型以及富金属的 M 型。

小行星带是由原始太阳星云中的一群微行星形成的$^{[9]}$，这些微行星是原行星的较小前身。然而，在火星与木星之间，来自木星的引力扰动阻碍了这些天体进一步聚积成一颗行星$^{[9][10]}$，并赋予它们过量的动能，从而在碰撞时将微行星以及大多数初生原行星击碎。其结果是，在太阳系历史最初的一亿年内，小行星带原始质量的 99.9\% 被消耗殆尽$^{[11]}$。其中一些碎片最终进入太阳系内部区域，导致与类地行星的陨石撞击。只要小行星绕太阳的公转周期与木星形成轨道共振，其轨道就会持续受到显著扰动；在这些轨道距离处会形成所谓的柯克伍德空隙，因为小行星被“扫”入其他轨道上$^{[12]}$。
\subsection{观测史}
\begin{figure}[ht]
\centering
\includegraphics[width=6cm]{./figures/afb56775c2129f99.png}
\caption{1596 年，约翰内斯·开普勒（Johannes Kepler）基于自己对于行星轨道比例关系的直觉，认为在火星与木星的轨道之间存在一颗尚未被发现的行星$^{[13]}$。} \label{fig_belt_3}
\end{figure}
1596 年，约翰内斯·开普勒（Johannes Kepler）在其著作《宇宙秘典》（Mysterium Cosmographicum）中写道：“在火星与木星之间，我放置一颗行星”，明确表达了他认为在那里应当存在一颗行星的预测$^{[14]}$。在分析第谷·布拉赫（Tycho Brahe）的观测数据时，开普勒认为火星与木星轨道之间的间隙过大，不符合他关于行星轨道应当分布位置的模型$^{[15]}$。

在 1766 年他匿名发表的对夏尔·博内（Charles Bonnet）《自然沉思录》（Contemplation de la Nature）的译本附注中$^{[16]}$，维滕贝格的天文学家约翰·丹尼尔·提丢斯（Johann Daniel Titius）$^{[17][18]}$ 指出行星布局中似乎存在某种规律，这一规律后来被称为提丢斯–波得定律（Titius–Bode Law）。如果从 0 开始构造一个数列，然后依次取 3、6、12、24、48 等数（每次加倍），再给每个数字加上 4 并除以 10，就可以得到与当时已知行星轨道半径（以天文单位计）惊人接近的数值——前提是必须承认在火星（对应数列的 12）与木星（48）轨道之间存在一颗“缺失的行星”（对应数列的 24）。在附注中，提丢斯写道：“难道造物主会把那片空间空出来吗？绝不会。”$^{[17]}$

当威廉·赫歇尔（William Herschel）于 1781 年发现天王星时，其轨道与这一定律高度吻合，使得一些天文学家断定，在火星和木星轨道之间必然存在一颗行星$^{[19]}$。
\begin{figure}[ht]
\centering
\includegraphics[width=6cm]{./figures/7ecae5c9889c484d.png}
\caption{朱塞佩·皮亚齐（Giuseppe Piazzi），即小行星带中最大天体谷神星（Ceres）的发现者：谷神星曾被视为一颗行星，但随后被重新分类为小行星，并自 2006 年起被归类为矮行星。} \label{fig_belt_4}
\end{figure}
1801 年 1 月 1 日，西西里巴勒莫大学天文学主席朱塞佩·皮亚齐（Giuseppe Piazzi）发现了一个微小的移动天体，其轨道半径与该规律所预测的数值完全一致。他将其命名为“谷神星”（Ceres），取自罗马的丰收女神、亦为西西里的守护神。皮亚齐起初认为它可能是一颗彗星，但其缺乏彗发的特征表明它更像是一颗行星$^{[20]}$。由此，上述规律成功预测了当时八颗行星（即水星、金星、地球、火星、谷神星、木星、土星和天王星）的轨道半长轴。

与谷神星的发现同时，应弗朗茨·克萨韦尔·冯·扎克（Franz Xaver von Zach）的邀请，24 名天文学家组成了一个非正式团体，被称为“天上警察”（celestial police），其明确目的在于寻找更多行星；他们重点在火星与木星之间搜索，因提丢斯–波得定律预言该区应存在一颗行星$^{[21][22]}$。

大约 15 个月后，“天上警察”成员海因里希·奥伯斯（Heinrich Olbers）在同一区域发现了第二个天体——智神星（Pallas）。但与其他已知行星不同，谷神星与智神星在望远镜最高放大倍率下仍然呈现为光点，而非解析为圆面。除了快速移动外，它们与恒星几乎无差别$^{[23]}$。

因此，1802 年，威廉·赫歇尔（William Herschel）建议将它们划入一个独立类别，并以希腊语 *asteroeides*（意为“类星的”）为来源，为其命名为“asteroids”（小行星）$^{[24][25]}$。在完成对谷神星与智神星的一系列观测后，他得出结论$^{[26]}$：

无论是“行星”还是“彗星”这一称谓，都不能恰当地应用于这两个天体……它们与小恒星极为相似，以至于几乎无法将其区分开来。基于它们这种类星的外观，我取其名，称之为“小行星”；但我保留更改其名称的权利，若日后发现一个能更好表达其本质的术语。

到 1807 年，进一步观测又在该区域发现了两颗新天体：婚神星（Juno）与灶神星（Vesta）$^{[23]}$。然而，拿破仑战争中利林塔尔（Lilienthal）天文台的焚毁使这一早期发现阶段被迫终结$^{[27][23]}$。

尽管赫歇尔已提出“小行星”一词，但在接下来的几十年里，人们仍普遍将这些天体称为行星$^{[16]}$，并按发现顺序为其编号：1 谷神星、2 智神星、3 婚神星、4 灶神星。然而 1845 年天文学家卡尔·路德维希·亨克（Karl Ludwig Hencke）发现了第五个天体（5 义神星 Astraea），紧接着更多天体以加速的速度被发现。将它们都算作行星变得愈发不切实际。最终，它们被从行星名单中移除（最早由亚历山大·冯·洪堡在 1850 年代初提出），而赫歇尔所造的“小行星”（asteroids）一词逐渐成为通用称谓$^{[16]}$。

1846 年海王星（Neptune）的发现使提丢斯–波得定律在科学界迅速失去信誉，因为其轨道与该定律所预测的位置毫不相近。直至今日，这一定律仍无科学解释，天文学界普遍认为它只是巧合$^{[28]}$。
\begin{figure}[ht]
\centering
\includegraphics[width=6cm]{./figures/922f458ab47f5d40.png}
\caption{951 号小行星加斯普拉（Gaspra）——首个被航天器成像的小行星，其图像来自伽利略号（Galileo）在 1991 年飞掠期间拍摄；图像的颜色经过了夸张处理。} \label{fig_belt_5}
\end{figure}
“asteroid belt”（小行星带）这一表达在 19 世纪 50 年代初开始被使用，但确切由谁最先创造此术语已难以考证。其最早的英文用例似乎出现在亚历山大·冯·洪堡（Alexander von Humboldt）的《宇宙》（*Cosmos*）1850 年的英文译本（由 Elise Otté 翻译）中$^{[29]}$：“[…] 以及流星在 11 月 13 日与 8 月 11 日附近的规律出现，这些流星很可能属于一条与地球轨道相交、以行星速度运动的小行星带的一部分。”另一处早期的用例出现在罗伯特·詹姆斯·曼（Robert James Mann）的《天国知识指南》（*A Guide to the Knowledge of the Heavens*）中$^{[30]}$：“小行星的轨道位于一条宽广的空间带中，延伸于[…] 的极限之间。”美国天文学家本杰明·皮尔斯（Benjamin Peirce）似乎采用并推广了这一术语$^{[31]}$。

到 1868 年年中，已定位的小行星超过 100 颗，而 1891 年马克斯·沃尔夫（Max Wolf）将天体摄影技术引入小行星搜寻后，发现速度显著加快$^{[32]}$。到 1921 年已发现 1,000 颗小行星$^{[33]}$，1981 年达到 10,000 颗$^{[34]}$，并在 2000 年突破 100,000 颗$^{[35]}$。现代小行星巡天系统现已使用自动化方法，以不断增长的数量探测新的小行星。

2014 年 1 月 22 日，欧洲航天局（ESA）的科学家首次明确报告在小行星带最大天体谷神星（Ceres）上探测到水汽$^{[36]}$。该探测由赫歇尔空间天文台（Herschel Space Observatory）利用其远红外能力完成$^{[37]}$。这一发现出乎意料，因为通常认为会“喷出气流与羽状物”的是彗星而非小行星。正如其中一位科学家所言：“彗星与小行星之间的界线正变得越来越模糊。”$^{[37]}$
\subsection{起源}
\begin{figure}[ht]
\centering
\includegraphics[width=6cm]{./figures/4dd59d36fca97ed2.png}
\caption{小行星带示意图：展示了轨道倾角与距离太阳的关系，其中小行星带核心区域的天体以红色标示，其他小行星以蓝色标示。} \label{fig_belt_6}
\end{figure}
\subsubsection{形成}

1802 年，在发现智神星（Pallas）后不久，奥伯斯（Olbers）向赫歇尔（Herschel）和卡尔·高斯（Carl Gauss）提出：谷神星（Ceres）与智神星可能是曾经存在于火星—木星区域的一颗巨大行星的碎片，该行星在数百万年前因内部爆炸或彗星撞击而解体$^{[38]}$；而敖德山天文学家 K. N. Savchenko 则提出：谷神星、智神星、婚神星（Juno）和灶神星（Vesta）并非爆炸行星的碎片，而是“逃逸的卫星”$^{[39]}$。然而，摧毁一颗行星所需的能量极其巨大，加之小行星带的总质量仅约为地球月球质量的 4\%$^{[3]}$，均不支持这些假说。此外，小行星之间显著的化学差异也难以用同一母体行星来解释$^{[40]}$。

关于小行星带形成的现代假说与太阳系总体的行星形成过程相关，即与长期存在的星云假说相似：一团由星际尘埃与气体构成的云在引力作用下坍缩，形成一个旋转的物质盘，随后凝聚形成太阳与行星$^{[41]}$。在太阳系历史的最初几百万年里，微粒之间的黏性碰撞使其逐渐聚集成更大的团块；当这些团块的质量达到一定阈值后，它们便能够通过引力吸引其他物体，成为微行星（planetesimals）。这种引力聚积过程最终导致行星的形成$^{[42]}$。

但位于未来小行星带所在区域的微行星受到了木星引力的强烈扰动$^{[43]}$。轨道共振发生在小行星的公转周期与木星轨道周期构成整数比的区域，从而将小行星扰动至其他轨道；火星与木星之间存在大量此类轨道共振。当木星在形成后向内迁移时，这些共振会扫过整个小行星带，使其中天体的动力学激发程度上升，并提高其相互之间的相对速度$^{[44]}$。在某些区域，碰撞的平均速度过高，使得微行星更容易在碰撞中破碎而非聚积$^{[45]}$，从而阻止了一颗行星的形成。结果，它们继续绕太阳运行，并偶尔发生碰撞$^{[43]}$。

在太阳系早期历史中，小行星经历了一定程度的熔融，使得其内部物质发生按质量分异。一些母体天体甚至可能经历过阶段性的爆炸性火山活动，形成岩浆海。然而，由于这些天体体积较小，相比大型行星，其熔融期极为短暂，并普遍在约 45 亿年前、即形成后的最初数千万年内结束$^{[46]}$。2007 年 8 月，一项对据信源自灶神星的南极陨石中锆石晶体的研究发现，灶神星（进而小行星带）形成速度可能非常快——在太阳系形成后的 1000 万年内就已形成$^{[47]}$。
\subsubsection{演化}
\begin{figure}[ht]
\centering
\includegraphics[width=6cm]{./figures/c1ce78ad56022f73.png}
\caption{主小行星带的大型小行星——4 号灶神星（Vesta）} \label{fig_belt_7}
\end{figure}
小行星并不是原始太阳系的“原始样本”。它们自形成以来经历了显著的演化，包括内部加热（在最初几千万年内）、撞击造成的表面熔融、辐射造成的空间风化，以及微陨石的持续轰击$^{[48][49][50][51]}$。尽管一些科学家将小行星称为残余微行星（residual planetesimals）$^{[52]}$，但也有科学家认为它们已是独特的类群$^{[53]}$。

当前的小行星带被认为仅含有原始小行星带质量的一小部分。计算机模拟表明，原始小行星带的质量可能与地球相当$^{[54]}$。主要由于引力扰动，大部分物质在形成后的约 100 万年内被从小行星带中清除，仅剩不到原始质量的 0.1\%$^{[43]}$。自形成以来，小行星带的大小分布保持相对稳定；主带小行星的典型尺寸并未出现显著增大或减小$^{[55]}$。

与木星存在 4:1 轨道共振的位置（轨道半径 2.06 天文单位，AU）可视为小行星带的内边界。木星的引力扰动会使位于该处的天体被推入不稳定轨道。大多数形成在此间隙以内的天体在太阳系早期历史中被火星（其远日点为 1.67 AU）吸积，或被火星的引力扰动驱逐出该区域$^{[56]}$。洪伽利亚小行星族（Hungaria asteroids）位于 4:1 共振内侧更靠近太阳的位置，但由于其轨道倾角较高，得以避免被扰乱$^{[57]}$。

在小行星带形成初期，距离太阳约 2.7 AU 的位置形成了一个“雪线”（snow line），其温度低于水的冰点。形成于该半径以外的微行星能够积聚冰$^{[58][59]}$。2006 年，小行星带外侧雪线以外被发现存在一类彗星群，它们可能为地球海洋提供了水源。根据某些模型，地球形成早期由火山逸出所释放的水量不足以形成海洋，因此需要外部来源，如彗星轰击$^{[60]}$。

小行星带外部包含一些可能在过去几百年间被“植入”该区域的含冰天体。其中一个例子是类希尔达彗星（quasi-Hilda comet）362P/(457175) 2008 GO98，它被认为可能是一颗曾经的半人马型天体（centaur），通过与木星的近距离接触而被送入外小行星带$^{[61]}$。
\subsection{特性}
\begin{figure}[ht]
\centering
\includegraphics[width=6cm]{./figures/8cba3638ed75e30c.png}
\caption{} \label{fig_belt_8}
\end{figure}